% solutions/prop-a3-horseshoe.tex
% Standalone proof dossier; it also lifts into main.tex as a subfile.
\documentclass[../main.tex]{subfiles}
\begin{document}

% === SOLUTION HEADER (proof provenance) =====================================
%   ledger-node : prop:a3-horseshoe
%   refines     : prop:a3-horseshoe
%   bounded_by  : obs:heavy-tail-no-classical
%   evidence_run: none
%   checked_by  : agent
%   author      : /root/a3
%   reviewer    : /root/cross_review
%   review      : research/reviews/2026-08-21-a3-independent-agent-audit.md
%   review level: independent agent; no human or Lean review is claimed
%   date        : 2026-08-21
% ============================================================================

\section*{Solution: the sharp weighted horseshoe Poincar\'e constant}
\label{sec:sol-prop-a3-horseshoe}

\begin{theorem}
\label{thm:sol-prop-a3-horseshoe}
Let $\mathrm{HS}_\tau$ be the marginal law generated by
$\theta\mid\lambda,\tau\sim N(0,\tau^2\lambda^2)$ and
$\lambda\sim C^+(0,1)$.  Its density is
\[
 h_\tau(\theta)
 =\frac{e^aE_1(a)}{\pi\sqrt{2\pi}\,\tau},
 \qquad
 a=\frac{\theta^2}{2\tau^2},
 \qquad
 E_1(a)=\int_a^\infty\frac{e^{-s}}s\,\dd s .
\]
Then the optimal constant in
\[
 \Var_{\mathrm{HS}_\tau}(f)
 \le C_{\mathrm{HS}}(\tau)
 \int_{\mathbb R}(\tau^2+\theta^2)|f'(\theta)|^2
       \dd\mathrm{HS}_\tau(\theta)
\]
is
\[
 C_{\mathrm{HS}}(\tau)=4
 \qquad(\tau>0).
\]
Moreover, the associated one-sided Hardy quantity
\[
 B_{\mathrm{HS}}(\tau)
 =\sup_{r>0}\mathrm{HS}_\tau([r,\infty))
   \int_0^r\frac{\dd t}{(\tau^2+t^2)h_\tau(t)}
\]
is finite and satisfies $B_{\mathrm{HS}}(\tau)=B_{\mathrm{HS}}(1)$.
\end{theorem}

\begin{proof}
The density follows by integrating out the half-Cauchy scale.  Indeed, with
$a=\theta^2/(2\tau^2)$ and the substitution $s=\lambda^{-2}$,
\[
 \int_0^\infty
 \frac{e^{-a/\lambda^2}}{\lambda(1+\lambda^2)}\,\dd\lambda
 =\frac12\int_0^\infty\frac{e^{-as}}{1+s}\,\dd s
 =\frac12e^aE_1(a),
\]
which gives the displayed normalization after multiplying the Gaussian and half-Cauchy
constants.  The change of variables $\theta=\tau x$ sends $\mathrm{HS}_\tau$ to
$\mathrm{HS}_1$ and sends the weighted energy exactly to
$\int(1+x^2)|F'(x)|^2\dd\mathrm{HS}_1(x)$, where $F(x)=f(\tau x)$.
It also leaves the displayed Hardy product unchanged.  It is therefore enough to take $\tau=1$.

Put
\[
 H(a)=e^aE_1(a)=\int_0^\infty\frac{e^{-t}}{a+t}\,\dd t,
 \quad N=a^2+5a+2,
 \quad Q=a(a^2+6a+6).
\]
The rational bound needed below has an exact positive remainder.  Pointwise in $t$,
\[
 N-\frac Q{a+t}
 =a(t-1)-t^2+6t-4+\frac{t(t^2-6t+6)}{a+t}.
\]
The two polynomial groups integrate to zero against $e^{-t}\dd t$.  Since
$\{e^{-t}t^3\}''=e^{-t}t(t^2-6t+6)$, two integrations by parts give
\begin{equation}
\label{eq:sol-a3-positive-remainder}
 N-QH(a)
 =2\int_0^\infty\frac{e^{-t}t^3}{(a+t)^3}\,\dd t>0.
\end{equation}
There is no endpoint term: $e^{-t}t^3$ and its derivative vanish at zero and infinity.
Consequently
\begin{equation}
\label{eq:sol-a3-rational-upper}
 H(a)<\frac{a^2+5a+2}{a(a^2+6a+6)}.
\end{equation}

Use the intrinsic coordinate $x=\sinh y$.  Up to normalization, the transformed density and
Dirichlet form are
\[
 q(y)=\cosh y\,H\!\left(\frac{\sinh^2y}{2}\right),
 \qquad
 \mathcal E(F)=\int_{\mathbb R}|F'(y)|^2q(y)\,\dd y.
\]
Write $W=-\log q$, and for $y>0$ set
$x=\sinh y$, $c=\cosh y$, and $a=x^2/2$.  The identity
$H'(a)=H(a)-1/a$ gives
\[
 W'(y)=xc\left(\frac1{aH(a)}-1\right)-\frac xc.
\]
Bound \eqref{eq:sol-a3-rational-upper}, followed by direct algebra with
$a=(c^2-1)/2$, gives
\[
 \frac1{aH(a)}
 \ge\frac{a^2+6a+6}{a^2+5a+2}
 \ge1+\frac1{c^2}+\frac1{c(c+1)},
\]
because the difference in the second inequality is
\[
 \frac{c^2(c-1)^2+5c^2+2c+1}
 {c^2(c+1)(c^4+8c^2-1)}>0.
\]
It follows that
\begin{equation}
\label{eq:sol-a3-drift}
 W'(y)\ge\frac{x}{c+1}=\tanh(y/2),
 \qquad y>0.
\end{equation}

Consider the half-line generator $\mathcal LF=F''-W'F'$ and the positive function
$g(y)=\sinh(y/2)$.  By \eqref{eq:sol-a3-drift},
\[
 -\frac{\mathcal Lg}{g}
 =-\frac14+\frac12W'(y)\coth(y/2)\ge\frac14.
\]
For a smooth compactly supported function $\varphi=gu$ with trace zero at the origin, the
ground-state identity is therefore
\begin{align*}
 \int_0^\infty\left(|\varphi'|^2-\frac14\varphi^2\right)q\,\dd y
 &=\int_0^\infty g^2|u'|^2q\,\dd y \\
 &\quad+\int_0^\infty
 \left(-\frac{\mathcal Lg}{g}-\frac14\right)\varphi^2q\,\dd y
 \ge0.
\end{align*}
The singularity at zero causes no missing boundary term.  As $y\downarrow0$,
$q(y)\asymp\log(1/y)$, so both $q$ and $q^{-1}$ are locally integrable.  Finite-energy
functions have a trace, since
\[
 |F(y)-F(0)|^2
 \le\left(\int_0^y|F'|^2q\right)
     \left(\int_0^yq^{-1}\right).
\]
For a smooth trace-zero test, $u=\varphi/g=O(1)$ and the boundary term is
$qgg'u^2=O(y\log(1/y))\to0$.  Smooth trace-zero tests form a core of the trace-zero half-line
form domain; closure thus yields
\begin{equation}
\label{eq:sol-a3-dirichlet-gap}
 \int_0^\infty h^2q\,\dd y
 \le4\int_0^\infty|h'|^2q\,\dd y
 \qquad(h(0)=0).
\end{equation}

This Dirichlet estimate controls both parities.  It directly controls the odd sector.  If $f$ is
in the even half-line sector and has zero $q$-mean, put $h=f-f(0)$.  Then $h(0)=0$,
$\mathcal E(h)=\mathcal E(f)$, and
\[
 \int_0^\infty h^2q
 =\int_0^\infty f^2q+f(0)^2\int_0^\infty q
 \ge\int_0^\infty f^2q.
\]
Applying \eqref{eq:sol-a3-dirichlet-gap} to $h$, and then using the orthogonal parity
decomposition, proves that the full weighted spectral gap is at least $1/4$.

It remains to prove sharpness.  The integral expansion of $H$ gives
\[
 H(a)=a^{-1}-a^{-2}+O(a^{-3}),
 \qquad
 q(y)=K e^{-y}\{1+O(e^{-2y})\}\qquad(y\to\infty).
\]
Choose a nonzero $\chi\in C_c^\infty(0,1)$ and set
$F_{R,L}(y)=e^{y/2}\chi((y-R)/L)$.  Its support is $[R,R+L]$.  Direct substitution and
$\int_0^1\chi\chi'=0$ give
\[
 \frac{\int|F_{R,L}'|^2q}
      {\int F_{R,L}^2q}
 =\frac14+O(L^{-2})+O(e^{-2R}).
\]
After normalizing $q$ to a probability density, Cauchy--Schwarz gives
$(\int F_{R,L}q)^2\le(\int F_{R,L}^2q)\,q([R,\infty))$.
Thus subtracting the mean changes no energy and a vanishing fraction of the squared norm as
$R\to\infty$.  Sending first $R$ and then $L$ to infinity proves that the full gap is at most
$1/4$.  Hence it equals $1/4$, and the optimal weighted Poincar\'e constant is $4$.

Finally, $H(a)\asymp\log(1/a)$ at zero and $H(a)\sim a^{-1}$ at infinity.  Hence
$h_1^{-1}$ is locally integrable at zero, while
$\mathrm{HS}_1([r,\infty))\asymp r^{-1}$ and
$\int^r\{(1+t^2)h_1(t)\}^{-1}\dd t\asymp r$ in the tail.  The Hardy product is bounded,
which completes the remaining assertion.
\end{proof}

\paragraph{Obstruction respected.}
The theorem is weighted by $\tau^2+\theta^2$ and makes no classical Poincar\'e, LSI, or
$T_2$ claim for the Cauchy-tailed horseshoe law.

\end{document}
